% ============================================================
% The Spectral Power of the Digit Function
% ============================================================

\documentclass[12pt]{amsart}

\usepackage{amsmath,amssymb,amsthm}
\usepackage{booktabs}
\usepackage{microtype}
\usepackage{url}

\newtheorem{theorem}{Theorem}
\newtheorem{lemma}[theorem]{Lemma}
\newtheorem{proposition}[theorem]{Proposition}
\newtheorem{corollary}[theorem]{Corollary}
\theoremstyle{definition}
\newtheorem{definition}[theorem]{Definition}
\theoremstyle{remark}
\newtheorem{remark}[theorem]{Remark}

\title{The Spectral Power of the Digit Function}

\author{Alexander S.\ Petty}

\email{alexander.petty@gmail.com}
\date{October 2021 (revised April 2026)}
\subjclass[2020]{11A63, 11B83, 42A16}

\begin{document}
\begin{abstract}
We derive an exact closed-form expression for the spectral
power $\Phi(k)$ of the digit partition induced by
$\delta(r) = \lfloor br/p \rfloor$ on
$\mathbb{Z}/p\mathbb{Z}$. For $k \ne 0$,
\[
\Phi(k) = \sum_{d=0}^{b-1}
\left|\frac{\sin(\pi k n_d / p)}{\sin(\pi k / p)}\right|^2,
\]
where $n_d$ is the size of the $d$-th digit bin. At $k = 0$,
$\Phi(0) = \sum_d n_d^2 = S(p,b)$, recovering the bin-sum
from~\cite{paper3}. Parseval's identity gives
$\sum_k \Phi(k) = p(p-1)$.

The spectral power captures the magnitude of the bin
Fourier coefficients but not their phase. Its zeroth mode
$\Phi(0) = S(p,b)$ governs the alignment limit
of~\cite{paper3}, and a prime $p$ is digit-partitioning
in base $b$ if and only if $\Phi(k)$ is constant in $k$.
The cyclic autocorrelation $R(\ell)$ used
in~\cite{paper7} to compute the eigenvalue spectrum
requires the phase information discarded by squaring,
and is treated separately in the sequel via a
two-variable cross-spectral function whose
diagonal recovers $\Phi$. The present paper identifies
the additive Fourier analysis of
$\mathbb{Z}/p\mathbb{Z}$ as the natural setting for the
digit-bin partition and derives every closed form that
follows from magnitude data alone.
\end{abstract}
\maketitle

% ============================================================
\section{Introduction}

The digit partition of a prime $p$ in base $b$ is the
collection of bins
$B_d = \{r : 1 \le r \le p{-}1,\ \lfloor br/p \rfloor = d\}$
for $d = 0, 1, \ldots, b{-}1$. The bin sizes $n_d = |B_d|$
are integer invariants of the pair $(p,b)$. Earlier papers in
this series have used the bin sizes through several scalar
combinations: the bin-sum $S(p,b) = \sum_d n_d^2$ that
governs the alignment limit~\cite{paper3}, the alignment
$\alpha$ and pairwise alignment $\sigma$ of the coherence
decomposition~\cite{paper5}, and the cyclic autocorrelation
$R(\ell)$ used in the spectral structure paper~\cite{paper7}
to compute the eigenvalues of the cross-alignment
matrix~\cite{paper6}. Each of these invariants is some
contraction of the bin partition. None of them is a function
of a single Fourier mode of the partition itself.

This paper identifies that single Fourier mode. We define
the \emph{spectral power} $\Phi(k)$ of the digit
partition as the sum of squared magnitudes of the bin
indicators' additive Fourier coefficients, prove that
$\Phi(k)$ has a closed form in terms of Dirichlet kernels
weighted by the bin sizes, and show that the existing
scalar invariants are values or contractions of $\Phi$.
The bin-sum $S(p,b)$ is $\Phi(0)$. The digit-partitioning
characterization $p \le b+1$ is equivalent to constancy of
$\Phi(k)$ in $k$. The cyclic autocorrelation $R(\ell)$ is
not determined by $\Phi$ alone (squaring discards the phase
information that $R$ depends on) but is naturally expressed
through a two-variable extension of $\Phi$, the
cross-spectral function, which is the subject of the
sequel paper.

The technical content of this paper is the observation
that the bins $B_d$ are contiguous integer intervals
(because $\delta(r) = \lfloor br/p \rfloor$ is monotone
in $r$), so each bin indicator's Fourier transform is a
geometric sum, and its magnitude is a Dirichlet kernel.
Once that is in place, the closed form for $\Phi(k)$
follows in one line from the definition. The recovery of
$S(p,b)$ as $\Phi(0)$ is a corollary of the same
identity, and Parseval's identity gives the total
spectral energy $\sum_k \Phi(k) = p(p-1)$ as a sanity
check.

% ============================================================
\section{The Bin Partition and Its Fourier Expansion}

\begin{definition}
For a prime $p$ not dividing $b$, the \emph{bin indicator} of
digit $d$ is the function
$\mathbf{1}_{B_d} : \{1, \ldots, p{-}1\} \to \{0,1\}$ defined by
\[
\mathbf{1}_{B_d}(r) =
\begin{cases}
1 & \text{if } \lfloor br/p \rfloor = d, \\
0 & \text{otherwise}.
\end{cases}
\]
The bin $B_d$ is a contiguous interval of integers in
$\{1, \ldots, p{-}1\}$ of length $n_d$.
\end{definition}

Since $B_d$ is a contiguous interval, the additive Fourier
transform of $\mathbf{1}_{B_d}$ is a geometric sum.

\begin{lemma}\label{lem:fourier-bin}
The Fourier coefficient of $\mathbf{1}_{B_d}$ at frequency $k$ is
\[
\widehat{\mathbf{1}}_{B_d}(k)
= \sum_{r \in B_d} e^{-2\pi i k r / p}.
\]
For $k \ne 0$, this equals a Dirichlet kernel evaluation:
\[
|\widehat{\mathbf{1}}_{B_d}(k)|
= \left|\frac{\sin(\pi k n_d / p)}{\sin(\pi k / p)}\right|.
\]
For $k = 0$, $|\widehat{\mathbf{1}}_{B_d}(0)| = n_d$.
\end{lemma}

\begin{proof}
The bin $B_d$ is an interval of $n_d$ consecutive integers starting
at some $a_d$. The sum is geometric:
\[
\widehat{\mathbf{1}}_{B_d}(k) = e^{-2\pi i k a_d / p}
\cdot \frac{1 - e^{-2\pi i k n_d / p}}{1 - e^{-2\pi i k / p}}.
\]
Taking absolute values cancels the phase and gives the stated
Dirichlet kernel form.
\end{proof}

% ============================================================
\section{The Spectral Power}

\begin{definition}
The \emph{spectral power} of the digit function at frequency $k$
is
\[
\Phi(k) = \sum_{d=0}^{b-1}
|\widehat{\mathbf{1}}_{B_d}(k)|^2.
\]
\end{definition}

\begin{theorem}\label{thm:spectral-power}
The spectral power has the exact closed form
\[
\Phi(k) = \sum_{d=0}^{b-1}
\left|
\frac{\sin(\pi k n_d / p)}{\sin(\pi k / p)}
\right|^2
\]
for $k \ne 0$, and $\Phi(0) = \sum_{d} n_d^2 = S(p,b)$.
\end{theorem}

\begin{proof}
Immediate from Lemma~\ref{lem:fourier-bin} and the definition.
\end{proof}

\begin{corollary}\label{cor:parseval}
$\sum_{k=0}^{p-1} \Phi(k) = p(p-1)$.
\end{corollary}

\begin{proof}
By Parseval's identity applied to each bin indicator:
$\sum_k |\widehat{\mathbf{1}}_{B_d}(k)|^2 = p \cdot n_d$.
Summing over $d$: $\sum_k \Phi(k) = p \sum_d n_d = p(p{-}1)$.
\end{proof}

% ============================================================
\section{Worked Example at $p = 13$ in Base $10$}\label{sec:example}

To make Theorem~\ref{thm:spectral-power} concrete, we
compute $\Phi$ explicitly at $p = 13$ in base $10$. The
bin sizes $n_d = |\{r \in \{1, \ldots, 12\} :
\lfloor 10r/13 \rfloor = d\}|$ are
\[
(n_0, n_1, n_2, n_3, n_4, n_5, n_6, n_7, n_8, n_9)
= (1, 1, 1, 2, 1, 1, 2, 1, 1, 1).
\]
Eight singleton bins and two doubled bins, summing to
$\sum_d n_d = 12 = p - 1$. The bin-sum is
\[
\Phi(0) = S(13, 10) = 8 \cdot 1 + 2 \cdot 4 = 16.
\]

For $k \ne 0$, the singleton bins contribute Dirichlet
kernel value $1$, and the doubled bins contribute
$|2 \cos(\pi k/13)|^2 = 4 \cos^2(\pi k/13)$. Summing,
\[
\Phi(k) = 8 + 8 \cos^2\!\left(\frac{\pi k}{13}\right),
\qquad k = 1, 2, \ldots, 12.
\]
The total spectral energy is
$\sum_{k=0}^{12} \Phi(k) = 16 + \sum_{k=1}^{12}
[8 + 8\cos^2(\pi k/13)] = 16 + 96 + 8 \cdot 11/2
= 16 + 96 + 44 = 156$, in agreement with
Corollary~\ref{cor:parseval} since
$p(p-1) = 13 \cdot 12 = 156$. The constant $11/2$ in
the cosine sum comes from the standard identity
$\sum_{k=1}^{p-1} \cos^2(\pi k/p) = (p{-}2)/2$.

% ============================================================
\section{Equivalent Characterization of Digit-Partitioning Primes}

The spectral power gives a clean spectral characterization
of the digit-partitioning property.

\begin{proposition}\label{prop:dp-flat}
A prime $p$ is digit-partitioning in base $b$
(equivalently, $p \le b{+}1$) if and only if
$\Phi(k)$ is constant in $k$ for $k \in \{0, 1, \ldots, p{-}1\}$.
\end{proposition}

\begin{proof}
If $p$ is digit-partitioning then $n_d \le 1$ for every $d$,
so each Dirichlet kernel
$|\sin(\pi k n_d/p)/\sin(\pi k/p)|$ is identically $1$ at
nonzero $k$, and $\Phi(k) = p{-}1$ for $k \ne 0$.
Also $\Phi(0) = \sum_d n_d^2 = p{-}1$ since each
nonzero $n_d$ equals $1$. So $\Phi$ is constant.

Conversely, suppose $\Phi$ is constant. Then $\Phi(0) =
\sum n_d^2$ equals the common value $\Phi(k)$ for any
nonzero $k$. By Corollary~\ref{cor:parseval}, the
constant value is $\Phi(k) = p - 1$ for all $k$.
So $\sum_d n_d^2 = p - 1 = \sum_d n_d$, which forces
$n_d \in \{0, 1\}$ for every $d$ (since $n_d^2 = n_d$
only at these values). Hence each nonempty bin has exactly
one element, so the bin partition is a digit partition,
which by~\cite{paper2} is equivalent to $p \le b{+}1$.
\end{proof}

% ============================================================
\section{The Position of $\Phi$ in the Chain}

The spectral power is the magnitude-only data of the bin
Fourier coefficients. To recover the cyclic autocorrelation
$R(\ell)$ used in~\cite{paper7} for the eigenvalue spectrum,
the phase information discarded by squaring is needed. The
``same digit'' indicator expands as
\[
\mathbf{1}[\delta(r) = \delta(s)]
= \sum_d \mathbf{1}_{B_d}(r)\, \mathbf{1}_{B_d}(s),
\]
and substituting $s = b^{\ell} r \bmod p$ followed by the
inverse Fourier expansion of each indicator gives
\[
R(\ell) = \frac{1}{p}
\sum_{k=0}^{p-1}
\sum_d \widehat{\mathbf{1}}_{B_d}(k)\,
\overline{\widehat{\mathbf{1}}_{B_d}(b^{-\ell}k)}.
\]
The inner sum at fixed $\ell$ involves cross-frequency
products
$\widehat{\mathbf{1}}_{B_d}(k)\,
\overline{\widehat{\mathbf{1}}_{B_d}(b^{-\ell}k)}$,
which are not captured by $|\widehat{\mathbf{1}}_{B_d}(k)|^2
= \Phi(k)$ alone. The two-variable cross-spectral function
$G(k,k') = \sum_d \widehat{\mathbf{1}}_{B_d}(k)\,
\widehat{\mathbf{1}}_{B_d}(k')$, of which $\Phi$ is the
antidiagonal $G(k,-k)$, is the natural object that
contains the phase data; it is introduced in the
sequel paper, where the closed-form expression for
$R(\ell)$ as a sum of $G$ along an orbit path is
proved.

% ============================================================
\section{Unification of Earlier Invariants}

Several scalar invariants of the alignment program are
contractions of $\Phi$. The alignment limit
of~\cite{paper3} is governed by the zeroth mode
$\Phi(0) = S(p,b)$, since the limiting alignment
$\mathcal{L}(p,b) = (p - 1 + S(p,b))/(p(p-1))$ in the
equidistributed case is a function of $\Phi(0)$ alone.
The coherence decomposition $\alpha = F + \sigma$
of~\cite{paper5} expresses $\alpha$, $\sigma$, and $F$
through $\Phi(0)$ together with the orbit structure of
$\langle b \rangle$ in $(\mathbb{Z}/p\mathbb{Z})^*$. The
digit-partitioning characterization, by
Proposition~\ref{prop:dp-flat}, is the statement that
$\Phi$ is constant in $k$. The eigenvalue spectrum
of~\cite{paper7} is the DFT of the cyclic autocorrelation
$R(\ell)$, which is not determined by $\Phi$ alone but
which is naturally expressed through the cross-spectral
function whose diagonal is $\Phi$.

In each case, $\Phi$ encodes the interaction between the
Euclidean geometry of the digit function (the bin sizes
$n_d$, which are integer counts of consecutive residues)
and the additive harmonic analysis of
$\mathbb{Z}/p\mathbb{Z}$ (the Dirichlet kernel at each
frequency). It is the simplest single object from which
the magnitude-side invariants of the program can be
read.

% ============================================================
\section{Remarks}

\subsection*{The Dirichlet kernel and the digit function}

The Dirichlet kernel $D_N(x) = \sin(N\pi x)/\sin(\pi x)$, a
classical object in harmonic analysis~\cite{zygmund,iwaniec}, appears
because the digit bins are contiguous intervals. This is a
consequence of the monotonicity of the digit function
$\delta(r) = \lfloor br/p \rfloor$: the function is
nondecreasing in $r$, and each preimage $B_d$ is a
contiguous interval (possibly empty). The spectral power
inherits the Dirichlet kernel's properties: concentration at low
frequencies for large bins, oscillatory behavior at high
frequencies.

\subsection*{Additive vs multiplicative characters}

The Legendre symbol (a multiplicative character) was tested in
the spectral structure paper~\cite{paper7} and found insufficient
to determine $R(\ell)$. The
spectral power uses additive characters instead, which succeed
because the bin partition is defined by an additive condition
(contiguous intervals in $\mathbb{Z}/p\mathbb{Z}$). The digit
function bridges multiplicative structure (the orbit of $b$) with
additive structure (the bin intervals), and the spectral power is
the natural invariant of this bridge.  The interplay between
additive and multiplicative characters in number-theoretic contexts
is treated systematically in~\cite{iwaniec}.

\subsection*{Toward $L$-functions}

The spectral power $\Phi(k)$ is defined on
$\mathbb{Z}/p\mathbb{Z}$ and is symmetric ($\Phi(k) =
\Phi(p{-}k)$). Its values at specific $k$ can be related to
Ramanujan sums~\cite{ramanujan} and Kloosterman-type expressions. Whether the
generating function $\sum_k \Phi(k) z^k$ or the Dirichlet series
$\sum_p \Phi_p(k) p^{-s}$ (summing over primes at a fixed
frequency $k$) exhibit analytic continuation or functional
equations is an open question.

% ============================================================
\begin{thebibliography}{9}

\bibitem{paper1}
A.~S.~Petty,
Why the golden ratio selects the prime three,
research note, 2020.

\bibitem{paper2}
A.~S.~Petty,
Digit-partitioning primes and the alignment formula,
research note, 2020.

\bibitem{paper3}
A.~S.~Petty,
The alignment limit for all primes,
research note, 2020.

\bibitem{paper5}
A.~S.~Petty,
The coherence decomposition,
research note, 2021.

\bibitem{paper6}
A.~S.~Petty,
The cross-alignment matrix,
research note, 2021.

\bibitem{paper7}
A.~S.~Petty,
The spectral structure of fractional fields,
research note, 2021.

\bibitem{zygmund}
A.~Zygmund,
\emph{Trigonometric Series}, vols.~I--II,
3rd ed., Cambridge University Press, 2002.

\bibitem{ramanujan}
S.~Ramanujan,
On certain trigonometrical sums and their applications in the
theory of numbers,
\emph{Trans.\ Cambridge Philos.\ Soc.} \textbf{22} (1918), 259--276.

\bibitem{iwaniec}
H.~Iwaniec and E.~Kowalski,
\emph{Analytic Number Theory},
AMS Colloquium Publications, vol.~53, 2004.

\end{thebibliography}

\end{document}
